\subsection{Discusión de resultados y limitaciones}
\label{subsec:discusion-resultados}

\subsubsection*{Interpretación respecto a la hipótesis y los objetivos}

La evidencia de la \autoref{sec:resultados} apoya la hipótesis de que sustituir
PD especializados por una política neuronal común es viable en condiciones
conocidas, pero no produce un único comportamiento arquitectónico homogéneo. En
ID, MLP, GRU y LSTM completan las misiones y se sitúan en el rango de RMSE del
PD representativo, con diferencias por familia. En OOD, la ventaja relativa de las
redes aparece en recombinaciones moderadas y en dos de las tres hélices
corregidas, mientras que lemniscatas extremas y
\texttt{helix\_tight\_aggressive} muestran que la generalización depende del
tipo de novedad geométrica y de la exigencia dinámica.

Los objetivos de comparar especialización clásica, imitación y memoria temporal
quedan cumplidos en el plano observacional: existe trazabilidad desde
\texttt{comparison\_all\_runs.csv} hasta cada \texttt{metrics.json}. No se
afirma superioridad neuronal global porque el éxito de misión OOD agregado
(66,3\,\%) es inferior al de \texttt{test} (99,5\,\%) y porque varios fallos
combinan error de seguimiento con protecciones o límites de actitud.

\subsubsection*{Dependencia del experto y del conjunto de datos}

Las redes predicen fuerza deseada del experto, no minimizan directamente el
error de posición. Por tanto, su techo de desempeño está acotado por la calidad
de las demostraciones filtradas y por la cobertura del banco outer-force. La
fidelidad supervisada favorable de GRU/LSTM no se traduce automáticamente en
menor RMSE en \texttt{circle} o \texttt{lissajous}, lo que es coherente con el
desajuste entre distribución de estados visitados en entrenamiento y régimen
dinámico en bucle cerrado.

\subsubsection*{Valor y límites del enfoque híbrido}

Mantener el lazo interno clásico conserva interpretabilidad aeroespacial y
protecciones compartidas. La \autoref{fig:res-protections-ood} muestra que, en OOD,
la seguridad observada no puede atribuirse a la red sin considerar el mixer, los
límites de fuerza y la saturación de actuadores. En OOD extremo, las
terminaciones por saturación persistente (MLP/GRU en
\texttt{lemniscate\_combined\_extreme}) indican que la red puede demandar
esfuerzos incompatibles con el actuador pese al lazo interno estable.

\subsubsection*{Papel de la arquitectura neuronal}

La tesis dominante del capítulo es que \emph{la arquitectura decide}. MLP
alcanza el menor RMSE en varias familias ID y en hélices OOD moderadas, pero
activa clipping de inclinación y puede disparar el error en lemniscatas con
viento y drag elevados. GRU ofrece el mejor compromiso supervisado y buen
comportamiento en varios OOD de \texttt{lissajous} y \texttt{composite}. LSTM
no supera sistemáticamente a GRU y exhibe el peor seguimiento en
\texttt{lemniscate\_fast\_center\_yaw} (\autoref{fig:res-trajectory-lemniscate-mlp-lstm}).
Estas diferencias se interpretan como sensibilidad distinta a historia temporal,
ruido y retardo, no como una jerarquía teórica fija entre arquitecturas.

\subsubsection*{Validez, limitaciones y amenazas}

El simulador 6DOF en mundo ENU con cuerpo FRD es fiel al alcance del TFG, pero
no sustituye validación con hardware real. La batería OOD consta de diez
escenarios: la corrección de \texttt{max\_duration\_s} eliminó un sesgo de
contrato temporal en dos hélices, pero \texttt{helix\_tight\_aggressive}
permanece como misión incompleta incluso con 60\,s, lo que limita las
conclusiones sobre waypoint agresivo bajo viento fuerte. Los RMSE agregados
muestran dispersión entre escenarios; no se ha cuantificado incertidumbre
estadística porque cada celda corresponde a una ejecución determinista por
escenario.

La comparación excluye \texttt{neural\_position} y no reentrena las redes tras
ajustar OOD: los checkpoints permanecen congelados. Amenazas adicionales
incluyen la elección de PD representativo por familia OOD
(\texttt{lemniscate}$\rightarrow$\texttt{lissajous},
\texttt{composite}$\rightarrow$\texttt{lissajous},
\texttt{waypoint}$\rightarrow$\texttt{waypoint}) y la dependencia de un único
par de semillas y hiperparámetros mínimos.

\subsubsection*{Implicaciones para la reproducibilidad}

La regeneración de \texttt{data/neural\_ood/battery\_v1}, la reejecución de
transferencias clásicas y de MLP/GRU/LSTM, y la consolidación con
\texttt{tools/summarize\_comparison.py} permiten reconstruir las tablas y figuras
citadas. Los comandos exactos y las rutas de artefactos se recogen en el
\apartadoref{anejo:comandos}. La separación entre evidencia
(\autoref{sec:resultados}) e interpretación (este apartado) busca que un lector
pueda auditar números sin mezclar observación con juicio causal.
